\documentclass[11pt]{article}
\usepackage[margin=1in]{geometry}
\usepackage[T1]{fontenc}
\usepackage{lmodern}
\usepackage{amsmath,amssymb}
\usepackage{enumitem}
\usepackage[hidelinks]{hyperref}
\usepackage{url}
\usepackage{xcolor}
\setlength{\emergencystretch}{3em}

%% Long Lean identifiers wrap only at their natural boundaries: after `.', `_'
%% and `-', and before `/'.  Allowing these breaks removes the need for \sloppy,
%% whose infinite tolerance produced very loose interword spacing.
\newcommand{\LeanBreakSetup}{%
  \urlstyle{tt}%
  \mathchardef\UrlBreakPenalty=100
  \mathchardef\UrlBigBreakPenalty=100
  \def\UrlBigBreaks{\do\/}%
  \def\UrlBreaks{\do\.\do\_\do\-}%
}
\DeclareUrlCommand\code{\LeanBreakSetup}

\newcommand{\review}[1]{\par\medskip\noindent
  \textbf{\color{blue!55!black}Reviewer comment.} \emph{#1}\par}
\newcommand{\response}{\par\smallskip\noindent
  \textbf{\color{green!35!black}Response.} }

\title{Response to Reviewer 1\\
\large ``The Seifert--van Kampen Theorem via Computational Paths''}
\author{Arthur F. Ramos, Tiago M. L. de Veras,\\
Ruy J. G. B. de Queiroz, and Anjolina G. de Oliveira}
\date{\today}

\begin{document}
\maketitle

We thank the reviewer for the careful and constructive report.  We retained the
paper's structure and addressed each identified location directly.  The Lean
development now includes proofs of both glue-naturality forms, the obstruction
to an amalgamation-only SVK target, the full-target theorem for the global
rewrite quotient, the low-degree \(\pi_0\) tail of the fibration sequence, and
a general sphere loop-quotient theorem.  It also includes the unconditional
space-level circle realization theorem described below.

As a final claim-level audit, we corrected two possible over-readings in the
extended-feature sections.  The revised manuscript restricts the
Eckmann--Hilton theorem to the explicitly defined \(\pi_2\), labels the
\(n\geq3\) branches of \code{HigherHomotopy.PiN} as current unit placeholders,
and presents the fibration construction as a low-degree sequence with
pointed adjacent-composition laws rather than a full image-kernel exactness
theorem.  We also moved the HoTT comparison table inside Section~10.1 and
made the formalization architecture explicitly Table~2.  For navigation, the
principal claim repairs are in Sections~2.5, 7.1, 7.5, 9.1, and 10.1, with the
declaration-level mapping in Appendix~A.

The revision also makes the quotient being classified explicit in every
nontrivial application.  The relation generated by all constructors of the
global \code{Path.Step} system makes each \code{PathRwQuot} loop fiber
contractible.

To obtain a non-collapsing fundamental-group theorem without relabeling a
normal-form quotient, we added \code{PresentedFundamentalGroup.lean}.  A
presented path space has proof-relevant generating edges; raw paths are freely
generated by reflexivity, reversal, and composition; and homotopy is generated
only afterward from named relators and the groupoid laws.  Its fundamental
group is the automorphism group of the basepoint in the resulting strict
groupoid.  The new theorem \code{presentedSeifertVanKampenGroupEquiv}
constructs encode and decode, proves both round trips, and records preservation
of multiplication, identity, and inversion:
\[
\pi_1^{\mathrm{pres}}(\mathcal P(G_1,G_2;H),*)
\cong G_1*_H G_2.
\]
We instantiate it for the figure-eight and separately prove a nontrivial
generator.  The circle presentation is likewise proved equivalent to
\(\mathbb Z\).  At the fundamental-group level, we additionally prove that
\(\mathbb R\to\mathsf{AddCircle}\,1\) is a covering map, classify Mathlib's
topological fundamental group of the circle by winding, and prove the group
equivalence \code{presentedCircleTopologicalPiOneGroupEquiv}.  A general
geometric realization is also constructed from the nerve, and
\code{Presented.Realization.hoNerveIso} recovers the original presented
groupoid as its simplicial homotopy category.  The final comparison with the
topological fundamental groupoid is now represented by an explicit canonical
functor, whose identity and composition laws, essential surjectivity, fullness,
and faithfulness are proved.  We descend degeneracy-aware open core-face stars
through the realization quotient, prove that they cover, partition each
preimage into disjoint lifted sheets, construct the sheet homeomorphisms, and
prove that the under-category projection is a Mathlib covering map.  Path and
homotopy lifting over its contractible total space then provide the explicit
hom-set inverse and the public equivalence
\code{topologicalFundamentalGroupoidEquivalence}.
This closes the general nerve-realization/fundamental-groupoid gap that the
previous draft listed under both Limitations and Future Work.  Separately,
\code{CircleNerveAmbient.lean} now proves the concrete ambient statement
\[
\mathsf{CircleNerveRealization}\simeq_h\mathsf{TopologicalCircle}.
\]
The maps \code{circleNerveToAmbient} and \code{circleAmbientToNerve} have an
exact right inverse, while a face- and degeneracy-compatible simplexwise join
homotopy descends through the realization colimit to prove the other inverse
up to homotopy.  Thus the circle case is no longer an ambient-space limitation;
the analogous comparison for arbitrary concrete CW presentations remains
Future Work.

We also completed the topological structure carried by the computational paths
themselves.  The new declaration
\code{unconditionalTopologicalGroupoidCertificate} packages continuous source,
target, identity, inverse, and composition together with the endpoint, unit,
inverse, and associativity laws on quotient classes.  It is specialized to all
presented realizations by
\code{continuousPresentedTopologicalGroupoidCertificate}, and transported
along any supplied ambient homotopy equivalence by
\code{continuousAmbientTopologicalGroupoidCertificate}.  The unconditional
composition domain has the final topology induced from composable
representatives.  We state separately that continuity for the ordinary
product/subspace topology retains the explicit
\code{ProductQuotientCompatibility} hypothesis, since products of quotient
maps are not quotient maps in arbitrary topological spaces.

Because these additions change where the paper places its weight, we also
added two short remarks that state their status plainly.  Remark~2.10 records
that the presented path space is the standard combinatorial presentation of a
groupoid, citing Higgins and Brown, so that the reader can see it is not an
object introduced to make a theorem succeed; Brown states Seifert--van Kampen
for fundamental groupoids in exactly this form.  Remark~2.7 gives the precise
mechanism of the global-rule collapse and reads it as a constraint on the
design space rather than as a defect of one implementation.

\section*{Overall remarks}

\review{The assumption bookkeeping is not applied consistently, and the reader
must consult several separated tables and appendices to determine what is
proved.}
\response
We consolidated the bookkeeping in four coordinated places: the main
presented theorem prints its algebraic data directly, the separate global-rule
schema prints all encode typeclasses in its statement, Section~8 contains one
manifest for those conditional parameters, and Appendix~A gives, for every
headline result, the exact Lean name, module, object classified, and status
(``presented fundamental group'', ``global-rule quotient'', ``conditional'',
``completed presentation'', or ``open'').
The text explains that the amalgamation-only relation preserves word length
and cites the formal impossibility theorems for that target.

\review{Claims about decidability and the ``36 axioms'' count are weaker than
the headline presentation suggests.}
\response
Both claims were corrected.  Arbitrary \(\mathsf{RwEq}\) is no longer described
as decidable without qualification: the paper now states that expression-level
termination and confluence are proved, while path-level decidability requires
an explicitly supplied \code{CompleteStepSystem}.  The ``36 axioms'' headline
and optional axiom-import inventory were removed.  The source contains
zero custom \code{axiom} declarations, as checked by
\code{scripts/check_invariants.py} and \code{paper/svk/check_stats.py}.

\review{A key naturality lemma is asserted without proof.}
\response
We added a complete derivation in the paper, together with the two Lean
theorems \code{Pushout.glue_natural_square_rweq} and
\code{Pushout.glue_natural_rearranged_rweq}.  Both are proved in the module
\code{Path/CompPath/PushoutCompPath.lean}.

\section*{Detailed comments}

\review{p.~2, Section 1, item 2: ``\(\sim\) is decidable via
normalization'' conflicts with the metatheory table.}
\response
Item~2 states that the normalization function is a sound invariant,
not an unconditional complete characterization.  It identifies the proved
abstract-expression fragment and points to
\code{CompleteStepSystem.decideRwEq} for conditional decidability.

\review{p.~10, Section 2.4: confluence is marked ``typeclass'', but it is not
stated who discharges it.}
\response
The table now distinguishes two levels.  The library discharges expression
termination and confluence through \code{instHasTerminationProp} and
\code{instHasConfluencePropExpr}.  A caller must supply
\code{HasJoinOfRw}/\code{HasConfluenceProp} at the trace-carrying
\code{Path} level; no global instance is claimed.

\review{p.~16, Lemma 3.4: no proof is given; add a proof or Lean pointer.}
\response
The added proof sets
\(L=\mathsf{inl}_*(f_*(p))\) and
\(R=\mathsf{glue}(c_1)\cdot\mathsf{inr}_*(g_*(p))\), uses the stored witness
\(L^{-1}\cdot R\sim\mathsf{glue}(c_2)\), and then applies left congruence,
associativity, inverse cancellation, and the unit law.  The paper cites
\code{glue_natural_square_rweq}.

\review{p.~16, Corollary 3.5: state explicitly that it follows by
right-multiplying by \(\mathsf{glue}(c_2)^{-1}\).}
\response
The proof now says exactly this, including reassociation and inverse
cancellation.  The corresponding Lean theorem is
\code{glue_natural_rearranged_rweq}.

\review{p.~21, Section 5.1: Theorem 5.1 gives no indication that it is
conditional.}
\response
The revised Section~5 now begins with an unconditional theorem for presented
computational path spaces.  Its raw paths and generated homotopy are defined
before encode, and Lean theorem
\code{presentedSeifertVanKampenGroupEquiv} constructs both maps, proves both
round trips, and packages the group homomorphism laws.  Its hypotheses
\code{GroupLaws} and \code{AmalgamationLaws} are explicit algebraic structures;
the latter includes the source-group laws and two
\code{PreservesGroupOps} certificates.  None of these are encode assumptions.
For the circle, the separate theorem
\code{presentedCircleTopologicalPiOneGroupEquiv} identifies the presented
fundamental group with Mathlib's topological
\texttt{FundamentalGroup (AddCircle 1) 0}.

The separate theorem for the global-rule \code{PathRwQuot} is titled
``Conditional global-rule SVK schema'', and all four required typeclasses are
part of its statement.  We also distinguish the full target, which includes
free-group reduction, from the amalgamation-only target.  Lean proves that
round-trip and bijectivity interfaces for the latter are impossible:
\code{hasPushoutSVKEncodeDecode_impossible} and
\code{hasPushoutSVKDecodeAmalgBijective_impossible}.  Separately,
\code{seifertVanKampenFullEquiv_collapsed} proves the full-target equivalence
for the global rewrite quotient and is explicitly labeled degenerate.
The displayed full-target definition now reproduces the six actual
\code{FreeGroupStep} constructors---adjacent combination on both sides, zero
removal on both sides, and left/right contextual congruence---and the explicit
reflexive, amalgamation, free-group, symmetric, and transitive constructors of
\code{FullAmalgEquiv}.
The completed-expression theorem \code{scopedSeifertVanKampenEquiv} remains as
a separate normal-form result rather than serving as the fundamental-group
theorem.

\review{p.~27, Table 1: the sphere row says ``SVK: Yes'', while Appendix A
lists no dependency for \(S^2\).}
\response
We added the general Lean theorem
\code{sphereN_piOne_equiv_unit}, specialized from
\code{pathRwQuotLoopEquivUnit}.  Table~1 and Appendix~A agree: this global-rule
loop-quotient result has no SVK encode/decode dependency and follows from the
universal quotient collapse.

\review{p.~36, Definition 7.6: the long exact sequence is truncated at
\(\to1\); should it continue through \(\pi_0\)?}
\response
Yes.  The display continues
\[
\pi_1(B)\to\pi_0(F)\to\pi_0(E)\to\pi_0(B).
\]
The revised subsection is now titled ``Low-degree fibration sequence''.  The
Lean development defines \code{boundaryPi0},
\code{fiberPi0ToTotalPi0}, and \code{totalPi0ToBasePi0}, and packages them with
the pointed adjacent-composition laws in
\code{LowDegreeFibrationSequence}; the constructor is
\code{lowDegreeFibrationSequence}.  It does not claim full image-kernel
exactness for the displayed sequence; that strengthening remains future work.
The paper explains that a terminal \(1\) is shorthand only under
path-connectedness assumptions (revised Section~7.5).

\review{p.~5, related work: add Agda/HoTT--Agda SVK developments and Lean
Mathlib's fundamental groupoid material.}
\response
Section~1.2 discusses the machine-checked \code{VanKampen} development in
HoTT--Agda and Mathlib's classical fundamental groupoid/fundamental group
modules, and explains how their objects differ from explicit equality-trace
quotients.

\review{No repository link is provided.}
\response
The abstract and a dedicated Code Availability subsection now cite the exact
Zenodo version~0.4.1 matching this revision:
\url{https://doi.org/10.5281/zenodo.21779002}.
It contains the 31 paper entry modules and their 61 transitive local
dependencies (92 Lean files total), including the topological computational
path groupoid package, \code{CircleNerveAmbient.lean}, the pinned toolchain and
Lake metadata, theorem manifest, license, invariant checker, and the corrected
manuscript and response sources.  Its archive SHA-256 is recorded in the
Zenodo metadata.  The concept
DOI for the archival series is
\url{https://doi.org/10.5281/zenodo.21696874}; the corresponding source release
is
\url{https://github.com/Arthur742Ramos/ComputationalPathsLean/releases/tag/v0.4.1},
and the repository is
\url{https://github.com/Arthur742Ramos/ComputationalPathsLean}.

\review{p.~44, Section 8.6 Part C: optional axioms are excluded from the
headline count and should be flagged next to each theorem.}
\response
The source contains no such optional axiom files, so Part~C and
the ``36 axioms'' count were removed rather than merely relocated.  Each
application is labeled locally as presented fundamental group,
global-rule/proved, conditional, completed presentation, or open, and
Appendix~A repeats the same status.

\review{p.~36, path tactics: state what fraction of theorems rely on them.}
\response
We added a reproducible lexical count.  None of 20,389
\code{theorem}/\code{lemma} declarations directly invokes a \code{path\_*}
tactic.  Across all 54,249 declarations counted, including definitions,
instances, and examples, 11 do so (approximately \(0.020\%\)); in the thirty-one
paper modules, one of 1,465 declarations does so (\(0.068\%\)).  The paper
therefore characterizes these tactics as convenience wrappers rather than a
mature dominant automation layer.

\review{The seventh line of the table on p.~51 does not express Theorem 6.8.}
\response
The Klein-bottle row and the corresponding theorem now name the declaration
\code{kleinBottleScopedPiOneEquivIntProd}, and identify its source as the
scoped expression quotient in \code{ClassicalPresentationsScoped.lean}.  Both places
state that the quotient is classified by the
\(\mathbb Z\rtimes\mathbb Z\) normal form and do not identify it with the
global-rule \code{PathRwQuot}.

\section*{Minor corrections}

\begin{itemize}[leftmargin=*]
  \item The structure paragraph now includes Section~9 before Section~10.
  \item The decode subsection now begins ``For pushouts, the decode map
        converts a word to a loop'', correcting the punctuation.
  \item The bullets in Example~6.12 begin on lines separate from the heading.
  \item ``Part C'' was removed with the optional-axiom
        inventory, eliminating the punctuation issue.
  \item Future Work now spells out the octonionic sequence
        \(S^7\to S^{15}\to S^8\).
\end{itemize}

\end{document}
